%% toc: P1 5 5.7
\section{Transfinite recursion on an ordinal: existence and uniqueness}
\label{sec:p01-w02-p01-013-transfinite-recursion}

\Cref{thm:transfinite-recursion} in \cref{sec:set-prerequisites-05} is the
principle by which later chapters define a function stage by stage along an
ordinal: the Lindenbaum and Henkin constructions of
\cref{ch:fol-metatheory} (\cref{lem:transfinite-lindenbaum}) are instances.
The proof is given in full below, together with the role of each hypothesis.
Two counterexamples show that unrestricted recursion along an ill-founded
order need not have either existence or uniqueness. These counterexamples
do not exclude recursion on well-founded relations that are not linear
orders.

\emph{Inputs.} An ordinal \(\delta\) and a functional rule \(G\) that assigns
a set \(G(g)\) to every function \(g\) whose domain is an ordinal smaller
than \(\delta\) (\cref{def:p01-w02-p01-013-reading}).
\emph{Conclusions.} Exactly one function \(F\) with domain \(\delta\)
satisfies \(F(\alpha)=G(F\mathbin{\upharpoonright}\alpha)\) for every
\(\alpha<\delta\) (\cref{thm:p01-w02-p01-013-recursion}); the functions
obtained for shorter lengths are restrictions of \(F\)
(\cref{thm:p01-w02-p01-013-coherence}); and a definition by separate zero,
successor and limit clauses is an instance
(\cref{thm:p01-w02-p01-013-cases}).

\subsection{Assumptions and vocabulary}
\label{sec:p01-w02-p01-013-assumptions}

\begin{hbconvention}{def:p01-w02-p01-013-reading}{Functions, restrictions and functional rules}
\begin{enumerate}
  \item The ambient theory is the classical first-order set theory of the
        metalanguage, in the language with the single binary relation symbol
        \(\in\). The proofs of the theorem, its lemmas and its corollaries use
        Extensionality, Pairing (to form ordered pairs), the Separation schema
        and the Replacement schema of \cref{sec:zf-zfc-01}, together with the
        definitions of transitive set and ordinal
        (\cref{def:transitive-set,def:ordinal-cardinal}). They use neither the
        Axiom of Choice nor Foundation nor Power Set. The worked example also
        uses Union and the set \(\N\) of \cref{def:p01-w02-p01-008-naturals,thm:p01-w02-p01-008-omega}.
  \item A \emph{function} is a set \(f\) of ordered pairs no two of which
        have the same first coordinate. Its domain \(\dom f\) is the class of
        first coordinates, and \(f(x)\) is the second coordinate of the pair
        with first coordinate \(x\). For a set \(A\), the \emph{restriction}
        \(f\mathbin{\upharpoonright}A\) is the set of pairs in \(f\) whose
        first coordinate lies in \(A\); it exists by Separation from \(f\).
  \item For ordinals, \(\alpha<\beta\) means \(\alpha\in\beta\)
        (\cref{def:ordinal-cardinal}), and \(\alpha\leq\beta\) means
        \(\alpha<\beta\) or \(\alpha=\beta\). The facts about ordinals used
        below are those of \cref{thm:p01-w02-p01-011-ordinals}.
  \item A \emph{functional rule} for an ordinal \(\delta\) is given by a
        formula \(\theta(g,y,\bar p)\) of the set-theoretic language, with
        set parameters \(\bar p\), such that for every function \(g\) whose
        domain is an ordinal smaller than \(\delta\) there is exactly one set
        \(y\) with \(\theta(g,y,\bar p)\). That set is written \(G(g)\), and
        \(G(g)=y\) abbreviates \(\theta(g,y,\bar p)\). A set function whose
        domain contains all such \(g\) is the special case in which
        \(\theta(g,y,G)\) says that the pair \((g,y)\) belongs to \(G\). The
        rule need not be defined on any other set, and it need not be given by
        a set.
\end{enumerate}
\end{hbconvention}

Because a functional rule is a formula, \cref{thm:p01-w02-p01-013-recursion}
below is a schema: one theorem of the ambient theory for each formula
\(\theta\). Every reference to \(G(g)\) in a proof unfolds to \(\theta\).

\begin{hbdefinition}{def:p01-w02-p01-013-attempt}{Attempt of a given length}
Let \(\delta\) be an ordinal and \(G\) a functional rule for \(\delta\). For
an ordinal \(\gamma\leq\delta\), an \emph{attempt of length \(\gamma\)} is a
function \(f\) with \(\dom f=\gamma\) such that
\[
  f(\alpha)=G(f\mathbin{\upharpoonright}\alpha)
  \qquad\text{for every }\alpha<\gamma.
\]
The right-hand side is defined: for \(\alpha<\gamma\), \(\alpha\) is an
ordinal smaller than \(\delta\) and \(f\mathbin{\upharpoonright}\alpha\) is a
function with domain \(\alpha\) (\cref{thm:p01-w02-p01-013-agreement},
item~1).
\end{hbdefinition}

An attempt of length \(\delta\) is exactly a function \(F\) as in the
conclusion of the theorem. The proof constructs attempts of every length
below \(\delta\) by a least-counterexample argument and then assembles them.

\subsection{The theorem}
\label{sec:p01-w02-p01-013-theorem}

\begin{hbtheorem}{thm:p01-w02-p01-013-recursion}{proved=annex:p01-w02-p01-013-recursion}{Transfinite recursion}
Let \(\delta\) be an ordinal. Suppose a functional rule \(G\) assigns a set
\(G(g)\) to every function \(g\) whose domain is an ordinal smaller than
\(\delta\). Then there is a unique function \(F\) with domain \(\delta\) such
that
\[
 F(\alpha)=G(F\mathbin{\upharpoonright}\alpha)
 \qquad\text{for every }\alpha<\delta.
\]
The theorem is the exact metalinguistic recursion principle used in later
chapters; it does not add a new inference rule to any object-language
calculus.
\end{hbtheorem}

The statement, its hypotheses and its notation are those of
\cref{thm:transfinite-recursion}. The proof rests on two lemmas.

\begin{hblemma}{thm:p01-w02-p01-013-agreement}{proved=annex:p01-w02-p01-013-agreement}{Restriction and agreement of attempts}
Let \(\delta\) be an ordinal and \(G\) a functional rule for \(\delta\).
\begin{enumerate}
  \item If \(\gamma\leq\delta\) and \(\alpha<\gamma\), then \(\alpha\) is an
        ordinal, \(\alpha<\delta\) and \(\alpha\subseteq\gamma\). If \(f\) is
        a function with domain \(\gamma\) and \(A\subseteq\gamma\), then
        \(f\mathbin{\upharpoonright}A\) is a function with domain \(A\) and
        \((f\mathbin{\upharpoonright}A)(x)=f(x)\) for \(x\in A\); and if
        \(B\subseteq A\), then
        \((f\mathbin{\upharpoonright}A)\mathbin{\upharpoonright}B
        =f\mathbin{\upharpoonright}B\).
  \item If \(f\) is an attempt of length \(\gamma\), \(f'\) is an attempt of
        length \(\gamma'\), and \(\alpha<\gamma\) and \(\alpha<\gamma'\), then
        \(f(\alpha)=f'(\alpha)\).
  \item For each \(\gamma\leq\delta\) there is at most one attempt of length
        \(\gamma\).
  \item If \(f\) is an attempt of length \(\gamma\) and \(\beta\leq\gamma\),
        then \(f\mathbin{\upharpoonright}\beta\) is an attempt of length
        \(\beta\).
\end{enumerate}
\end{hblemma}

\begin{hblemma}{thm:p01-w02-p01-013-assembly}{proved=annex:p01-w02-p01-013-assembly}{Assembling the earlier attempts}
Let \(\delta\) be an ordinal, \(G\) a functional rule for \(\delta\), and
\(\gamma\leq\delta\). Suppose that for every \(\alpha<\gamma\) there is an
attempt of length \(\alpha\), and write \(f_\alpha\) for it, which is unique
by \cref{thm:p01-w02-p01-013-agreement}. Then
\[
  h_\gamma=\{(\alpha,G(f_\alpha)):\alpha<\gamma\}
\]
is a set, and it is an attempt of length \(\gamma\). Moreover
\(h_\gamma\mathbin{\upharpoonright}\alpha=f_\alpha\) for every
\(\alpha<\gamma\).
\end{hblemma}

The Replacement schema is used once in the whole argument, to show that
\(h_\gamma\) is a set. The well-ordering of \(\delta\) by \(\in\) is used in
two places: to find a least point of disagreement in
\cref{thm:p01-w02-p01-013-agreement}, item~2, and to find a least length
without an attempt in the proof of the theorem.

\begin{hbcorollary}{thm:p01-w02-p01-013-coherence}{proved=annex:p01-w02-p01-013-coherence}{Shorter lengths give restrictions}
Let \(\delta\), \(G\) and \(F\) be as in
\cref{thm:p01-w02-p01-013-recursion}, and let \(\gamma\leq\delta\). Then
\(G\) assigns a set to every function whose domain is an ordinal smaller than
\(\gamma\), and the unique function \(F_\gamma\) that the theorem gives for
the ordinal \(\gamma\) and the same rule \(G\) is
\(F\mathbin{\upharpoonright}\gamma\).
\end{hbcorollary}

\begin{hbcorollary}{thm:p01-w02-p01-013-cases}{proved=annex:p01-w02-p01-013-cases}{Recursion by zero, successor and limit clauses}
Let \(\delta\) be an ordinal and \(a\) a set. Let \(S\) be a functional rule
that assigns a set \(S(\eta,x)\) to every ordinal \(\eta<\delta\) and every
set \(x\), and let \(L\) be a functional rule that assigns a set \(L(g)\) to
every function \(g\) whose domain is a limit ordinal smaller than \(\delta\).
Then there is a unique function \(F\) with domain \(\delta\) such that
\begin{enumerate}
  \item \(F(0)=a\) if \(0<\delta\);
  \item \(F(\eta+1)=S(\eta,F(\eta))\) for every ordinal \(\eta\) with
        \(\eta+1<\delta\);
  \item \(F(\lambda)=L(F\mathbin{\upharpoonright}\lambda)\) for every limit
        ordinal \(\lambda<\delta\).
\end{enumerate}
\end{hbcorollary}

The definitions in \cref{lem:transfinite-lindenbaum} and the ordinal
operations of \cref{def:p06-ordinal-arithmetic} are written in this
three-clause form.

\begin{hbcorollary}{thm:p01-w02-p01-013-partial}{proved=annex:p01-w02-p01-013-partial}{Recursion for a partial rule}
Let \(\delta\) be an ordinal. Let \(\theta(g,y,\bar p)\) be a formula that
has at most one value \(y\) for each function \(g\) whose domain is an
ordinal smaller than \(\delta\). For \(\gamma\leq\delta\), call a function
\(f\) with domain \(\gamma\) a partial-rule attempt if
\[
 \theta(f\mathbin{\upharpoonright}\beta,f(\beta),\bar p)
 \qquad(\beta<\gamma).
\]
The following conditions are equivalent:
\begin{enumerate}
 \item There is a function \(F\) with domain \(\delta\) satisfying
       \(\theta(F\mathbin{\upharpoonright}\alpha,F(\alpha),\bar p)\)
       for every \(\alpha<\delta\).
 \item For every \(\alpha<\delta\) and every partial-rule attempt \(f\)
       of length \(\alpha\), there is a set \(y\) with
       \(\theta(f,y,\bar p)\).
\end{enumerate}
When they hold, the function \(F\) is unique. Thus a value is required on
every history satisfying the earlier recursion equations, not necessarily
on every function of the same length.
\end{hbcorollary}

\subsection{Defining a function by recursion: procedure and side conditions}
\label{sec:p01-w02-p01-013-procedure}

\begin{hbmethod}{def:p01-w02-p01-013-procedure}{Defining a function along an ordinal}
\begin{enumerate}
  \item Fix the length \(\delta\) and check that it is an ordinal.
  \item Write the rule \(G\) as a formula \(\theta(g,y,\bar p)\) and check
        that it yields exactly one \(y\) for \emph{every} function \(g\)
        whose domain is an ordinal smaller than \(\delta\), not only for the
        restrictions of the function being defined. For a definition by
        cases, use \cref{thm:p01-w02-p01-013-cases} and check the three
        clauses separately.
  \item Conclude by \cref{thm:p01-w02-p01-013-recursion} that exactly one
        \(F\) with domain \(\delta\) satisfies the recursion equation.
  \item To prove a property of every value \(F(\alpha)\), consider the set
        of \(\alpha<\delta\) at which it fails, formed by Separation from
        \(\delta\), and refute its least member.
\end{enumerate}
The rule may inspect the entire earlier function, recover the current
stage from its domain, and use the fixed parameters \(\bar p\). Its formula
must determine a unique value from those data. This permits a rule whose
behavior varies with the earlier values; it does not permit dependence on
an unspecified later value. A requirement to select an arbitrary member of
a nonempty candidate collection is not yet a functional rule. One must
specify and justify a unique selection, for example by a definable
criterion or by a suitable choice function supplied as a parameter.
A well-order can provide such a criterion when every candidate collection
has a least member in that order. The recursion theorem alone supplies
neither the selector nor a well-order of all sets.
\end{hbmethod}

\subsection{Worked example: generating closure under successor}
\label{sec:p01-w02-p01-013-example}

\begin{hbexample}{def:p01-w02-p01-013-closure-example}{Stages of the closure of \(\{0\}\) under successor}
Natural numbers are the sets \(n=\{0,\ldots,n-1\}\) of
\cref{def:p01-w02-p01-008-naturals,thm:p01-w02-p01-008-omega}, with \(n+1=n\cup\{n\}\). For
\(X\subseteq\N\) put
\[
  \Phi(X)=\{0\}\cup\{k+1:k\in X\}.
\]
The map \(\Phi\) performs one generation step; it is not itself a closure
operator on the subsets of \(\N\). It is monotone: if \(X\subseteq Y\),
then every successor contributed by \(X\) is also contributed by \(Y\),
and both images contain \(0\). However, it need not be extensive:
\(\Phi(\{1\})=\{0,2\}\) does not contain \(1\). It is not idempotent either:
\(\Phi(\varnothing)=\{0\}\), whereas
\(\Phi(\Phi(\varnothing))=\{0,1\}\).
Here a closure operator means an extensive, monotone and idempotent map;
the closure generated by the rules must be distinguished from their
one-step application.
A subset \(X\) of \(\N\) with \(\Phi(X)\subseteq X\) contains \(0\) and is
closed under \(k\mapsto k+1\), so \(X=\N\) because \(\N\) is the least such
set. Thus \(\N\) is the only subset of \(\N\) closed under \(\Phi\).
Consequently, the least subset of \(\N\) containing any prescribed
\(X\subseteq\N\), containing \(0\), and closed under successor is \(\N\)
itself. The operation assigning this least closed superset is the actual
closure operator in this example. Its extensivity follows from
\(X\subseteq\N\), its monotonicity from its constant value, and its
idempotence from assigning \(\N\) to \(\N\).

Take \(\delta=(\N+1)+1\), the ordinal whose members are the natural numbers,
\(\N\) and \(\N+1\), and the rule
\[
  G(g)=\Phi\Bigl(\N\cap\bigcup\{g(\beta):\beta\in\dom g\}\Bigr),
\]
which assigns a set to every function \(g\). The theorem gives a unique
\(F\) with domain \(\delta\), and \(F(\alpha)\) is \(\Phi\) applied to the
union of all earlier values. Its values are
(\cref{thm:p01-w02-p01-013-closure-stages}):
\[
  \begin{array}{l|llll}
    \alpha & 0 & 1 & n & \N \text{ and } \N+1\\ \hline
    F(\alpha) & \{0\} & \{0,1\} & n+1=\{0,\ldots,n\} & \N
  \end{array}
\]
At each finite stage one new number is added, so no \(F(n)\) is closed under
\(\Phi\). Stage \(\N\) is the first limit stage. The same uniform clause
that produced the finite stages now applies \(\Phi\) to the union of all of
them, which is \(\N\), and the value is the closed set \(\N\). Stage
\(\N+1\) repeats it. No separate limit clause was written: the rule \(G\)
treats every stage alike.
\end{hbexample}

\begin{hblemma}{thm:p01-w02-p01-013-closure-stages}{proved=annex:p01-w02-p01-013-closure-stages}{Values in the closure example}
In \cref{def:p01-w02-p01-013-closure-example}, \(\N\) is a limit ordinal,
\(\N+1\) and \((\N+1)+1\) are ordinals, \(G\) is a functional rule for
\((\N+1)+1\), and the function \(F\) satisfies \(F(n)=n+1\) for every
\(n\in\N\), \(F(\N)=\N\) and \(F(\N+1)=\N\). Moreover
\(\Phi(F(n))\ne F(n)\) for every \(n\in\N\), and \(\Phi(F(\N))=F(\N)\).
\end{hblemma}

\subsection{Limitations}
\label{sec:p01-w02-p01-013-limits}

The proof finds a least position at which something fails. The next lemma
shows that both conclusions of the theorem can fail when the positions are
ordered without least elements. Index \(\N\) by the converse order: the
positions earlier than \(n\) are the members of
\(T_n=\{m\in\N:n<m\}\). The set \(\N\) itself then has no earliest
position, since \(n+1\in T_n\) for every \(n\).

\begin{hblemma}{thm:p01-w02-p01-013-ill-founded}{proved=annex:p01-w02-p01-013-ill-founded}{Ill-founded recursion}
For a function \(g\), let \(V(g)=\{g(x):x\in\dom g\}\) be its set of values.
Let \(G_1(g)=1\) if \(1\in V(g)\) and \(G_1(g)=0\) otherwise, and let
\(G_2(g)=0\) if \(1\in V(g)\) and \(G_2(g)=1\) otherwise.
\begin{enumerate}
  \item \emph{Uniqueness fails.} The constant function with domain \(\N\)
        and value \(0\) and the constant function with domain \(\N\) and
        value \(1\) both satisfy
        \(F(n)=G_1(F\mathbin{\upharpoonright}T_n)\) for every \(n\in\N\).
  \item \emph{Existence fails.} No function \(F\) with domain \(\N\)
        satisfies \(F(n)=G_2(F\mathbin{\upharpoonright}T_n)\) for every
        \(n\in\N\).
\end{enumerate}
\end{hblemma}

\begin{hbwarning}{def:p01-w02-p01-013-limits}{What the recursion theorem does not provide}
\begin{enumerate}
  \item \emph{Well-ordering is the indexing hypothesis here.} The theorem is stated for
        an ordinal, ordered by \(\in\). By
        \cref{thm:p01-w02-p01-013-ill-founded}, along an order with no least
        element in some nonempty set, existence and uniqueness can each fail.
  \item \emph{Totality is a sufficient hypothesis.} The theorem requires a
        value on every function of the indicated ordinal lengths. A partial
        rule has a unique solution exactly under the extension condition of
        \cref{thm:p01-w02-p01-013-partial}. Undefined values on histories
        that fail earlier recursion equations do not prevent that result.
        If \(G\) is
        undefined at some function whose domain is an ordinal below
        \(\delta\), the conclusion can fail. For \(\delta=2\), a rule with
        \(G(\varnothing)=\varnothing\) that is undefined at
        \(\{(0,\varnothing)\}\) admits no \(F\): the equation at \(0\) forces
        \(F(0)=\varnothing\), so \(F\mathbin{\upharpoonright}1=\{(0,\varnothing)\}\),
        and the equation at \(1\) has no right-hand side.
  \item \emph{One length at a time.} For each ordinal \(\delta\) the theorem
        produces a set function with domain \(\delta\). A function on the
        class of all ordinals is a different object; its statement is
        \cref{thm:p06-transfinite-recursion}, and
        \cref{thm:p01-w02-p01-013-coherence} supplies the agreement between
        lengths on which it depends.
  \item \emph{Existence, not computation.} The theorem does not say that the
        values are computable or explicitly describable; that depends on
        \(G\) and on how \(\delta\) is presented.
  \item \emph{Replacement.} The proof uses Replacement exactly once, to form
        \(h_\gamma\) in \cref{thm:p01-w02-p01-013-assembly}. Nothing proved
        in this section shows that this use can or cannot be avoided.
\end{enumerate}
\end{hbwarning}

\subsection{Relationships}
\label{sec:p01-w02-p01-013-relationships}

The theorem proved here is \cref{thm:transfinite-recursion}, with the same
hypotheses and conclusion. Its users include the
transfinite Lindenbaum construction (\cref{lem:transfinite-lindenbaum}) and
the Henkin recursion of \cref{ch:fol-metatheory}. The inductively generated
sets of \cref{sec:strings-02} are built by finite stages; the example above
shows how the same stage construction continues to a limit stage. In Part~6,
transfinite induction is \cref{thm:p06-transfinite-induction}, the
class-length recursion theorem is \cref{thm:p06-transfinite-recursion}, and
ordinal arithmetic (\cref{def:p06-ordinal-arithmetic}) is defined by the
three-clause form of \cref{thm:p01-w02-p01-013-cases}. The ordinal facts used
in the proofs are \cref{thm:p01-w02-p01-011-ordinals}.

\subsection{Sources}
\label{sec:p01-w02-p01-013-sources}

The statement of \cref{thm:p01-w02-p01-013-recursion}, with its hypotheses,
notation and closing qualification, is the inherited theorem
\cite[ch.~``Sets, Relations, Functions, and Cardinality'', sec.~``Finite,
countable, and uncountable sets'', theorem ``Transfinite
recursion'']{FrozenFormalLogicTransfiniteRecursion}. There it carries the
standing of a cited theorem and its proof is referred to \cite{Jech2003}.
The Separation and Replacement schemas used in the proofs are those of the
same manuscript \cite[ch.~``The Axioms of ZF and ZFC'', sec.~``The language
and the axiom schemas'', definitions ``Separation schema'' and
``Replacement schema'']{FrozenFormalLogicTransfiniteRecursion}. The reading of
``functional rule'' as a formula, the notion of an attempt, the proofs, the
two corollaries, the example and the counterexamples are editorial and are
not attributed to the manuscript.
